\section{Преглед на съществуващите решения}
\label{sec:existing_solutions}

Анализът на съществуващите решения е необходим, за да се оцени дали предложената архитектура действително допринася с нова стойност, или само преформулира вече познати практики. При C++ библиотеките за достъп до данни могат да се разграничат две основни оси и една допълнителна изпълнителна перспектива: нивото на абстракция, моментът, в който се изгражда и валидира заявката, и моделът на изпълнение. Нивото на абстракция варира от почти директно използване на клиентски C програмен интерфейс (API) до пълно поведение за обектно-релационно съпоставяне (ОРС). Моментът на валидиране може да бъде по време на изпълнение, хибриден или по време на компилация. Изпълнителният модел варира от директни блокиращи извиквания на програмния интерфейс до асинхронни автомати с крайни състояния, специфични за свързващия компонент, и по-високи слоеве, които се опитват да ги уеднаквят.

Тези оси са особено важни в контекста на настоящата дипломна работа. Ако се наблюдава само ``удобството на програмния интерфейс'', лесно може да се пропусне съществената разлика между библиотека, която просто прави SQL низовете по-четими, и библиотека, която действително прехвърля структурната коректност към компилатора. Ако, от друга страна, се гледа само кога възниква грешката, може да се подцени цената на различните нива на абстракция, сложността на инструментариума (инструментариум) за изграждане и пригодността за работа с повече от един модел на съхранение. А ако се игнорира изпълнителният модел, лесно се стига до погрешното впечатление, че ``асинхронност'' е еднотипна способност, независимо дали става дума за платформено-специфичен неблокиращ програмен интерфейс, слой с функции за обратно извикване върху цикъл за събития (цикъл за събития) или просто за изнасяне на работата към набор от нишки.

\subsection{Класификация на подходите}

В най-ниския слой стоят платформено-специфичните клиентски библиотеки като \code{SQLite C API}~\cite{sqlite_api}, \code{libpq}~\cite{libpq} и \code{MySQL Connector/C}~\cite{mysql_connector}. Те предоставят максимален контрол върху заявките, параметрите и резултатите, но не предлагат логическа абстракция. Приложението е принудено да работи директно със SQL низове, индекси на колони, позиции на параметри и типове, специфични за свързващия компонент.

Следващата група са библиотеките за конструиране на заявки по време на изпълнение. Те обикновено подобряват четимостта на заявките и намаляват част от шаблонния код, но продължават да разчитат на изграждане на текстово представяне по време на изпълнение. Типовата безопасност е частична и обикновено спира до границата на извикването на програмния интерфейс. Структурната валидност на заявката остава проблем на времето за изпълнение.

Най-високото ниво на абстракция е подходът за обектно-релационно съпоставяне (ОРС). Там C++ класовете се асоциират с таблици, колони и връзки, а библиотеката генерира необходимите заявки. Предимството е намаляване на шаблонния код. Недостатъкът е, че често се въвеждат външни генератори, макроси, анотации и непрозрачни метаданни по време на изпълнение, които затрудняват проследимостта между C++ кода и реалната заявка.

\begin{figure}[H]
\centering
\begin{tikzpicture}[x=2.0cm,y=1cm]
    \draw[->, thick] (0,0) -- (9,0) node[right] {\small момент на валидация};
    \draw[->, thick] (0,0) -- (0,7) node[above] {\small ниво на абстракция};

    \node[below] at (1.6,0) {\small по време на изпълнение};
    \node[below] at (4.6,0) {\small хибриден};
    \node[below] at (7.6,0) {\small по време на компилация};

    \node[left] at (0,1.2) {\small ниско};
    \node[left] at (0,3.6) {\small средно};
    \node[left] at (0,6.0) {\small високо};

    \node[ormblock, minimum width=2.2cm] at (1.8,1.4) {Специфичен\\програмен интерфейс};
    \node[ormblock, minimum width=2.6cm] at (3.6,3.4) {Конструктори\\на заявки};
    \node[ormblock, minimum width=2.4cm] at (5.6,4.2) {Типизирани SQL\\решения};
    \node[ormblock, minimum width=2.4cm] at (4.8,5.6) {Класическо ОРС\\+ кодогенерация};
    \node[ormaccent, minimum width=2.5cm] at (7.6,5.8) {ОРС по време на\\компилация};
\end{tikzpicture}
\caption{Позициониране на основните семейства решения според абстракция и момент на валидация}
\label{fig:existing_solutions_map}
\end{figure}

Фигура~\ref{fig:existing_solutions_map} показва, че подходът за обектно-релационно съпоставяне по време на компилация не е просто ``още една ОРС библиотека''. Той стои в горния десен квадрант, където високото ниво на абстракция се комбинира с ранна структурна проверка. Тази комбинация е рядка, защото изисква повече архитектурна дисциплина в самата библиотека.

\begin{table}[H]
\centering
\caption{Съпоставка на основните подходи за достъп до данни в C++}
\label{tab:existing_solutions_comparison}
\begin{tabular}{L{3.2cm}L{2.5cm}L{2.7cm}L{2.9cm}L{2.8cm}}
\toprule
\textbf{Подход} & \textbf{Ниво на абстракция} & \textbf{Валидация} & \textbf{Предимство} & \textbf{Основен недостатък} \\
\midrule
Специфичен програмен интерфейс & ниско & по време на изпълнение & пълен контрол и близост до драйвера & силна обвързаност с конкретен свързващ компонент \\
Конструктор на заявки & средно & основно по време на изпълнение & по-четим интерфейс и по-малко шаблонен код & заявката пак се материализира като низ \\
Типизирано SQL решение & средно към високо & частично по време на компилация & по-добра статична структура за SQL света & обикновено остава ограничено до SQL \\
Класическо ОРС & високо & по време на изпълнение или чрез кодогенерация & удобно съпоставяне на обекти & сложна интеграция в инструментариума и непрозрачност \\
ОРС по време на компилация & високо & по време на компилация за структурата & ранно откриване на структурни грешки и формален договор на конектора & по-висока сложност на шаблонния слой \\
\bottomrule
\end{tabular}
\end{table}

\subsection{Критерии за сравнение}

За да бъде сравнението академично пълно, не е достатъчно да се каже, че една библиотека е ``по-удобна'', а друга е ``по-ниско ниво''. В настоящата работа се използват шест сравнителни критерия:

\begin{enumerate}[label=\arabic*.]
    \item \textbf{Ниво на логическа абстракция} --- доколко библиотеката позволява моделът на данните да се описва независимо от непосредствения синтаксис на свързващия компонент.
    \item \textbf{Степен на статична валидация} --- кои аспекти на заявката и на модела могат да бъдат проверени преди изпълнение.
    \item \textbf{Преносимост между свързващи компоненти} --- възможно ли е логическият модел да бъде адаптиран към повече от един тип система за съхранение.
    \item \textbf{Сложност на инструментариума} --- изисква ли библиотеката външна кодогенерация, допълнителни стъпки за предварителна обработка (предварителна обработка) или генератори на метаданни.
    \item \textbf{Прозрачност на грешките и на изпълнението} --- доколко разработчикът може да проследи връзката между C++ кода, генерираната заявка и поведението по време на изпълнение.
    \item \textbf{Изпълнителен модел и асинхронна пригодност} --- дали библиотеката предоставя единен и проверим модел за конкурентно или неблокиращо изпълнение, или оставя тази задача изцяло на конкретния драйвер и приложния код.
\end{enumerate}

Тези критерии са пряко свързани с целите на дипломната работа. Ако предложената библиотека трябва да бъде едновременно изразителна и формално проверима, тогава тя трябва да превъзхожда алтернативите не само по удобство, а по баланса между тези шест измерения.

\subsection{Платформено-специфични клиентски библиотеки}

SQLite~\cite{sqlite}, \code{libpq} и \code{MySQL Connector/C} са представителни за подхода, при който приложението конструира и подава заявките директно към драйвера. Това е най-универсалният и често най-бързият вариант, но поставя тежестта на коректността изцяло върху приложния код. Разработчикът трябва ръчно да поддържа синхрон между имената на полета, поредността на параметрите и реалната схема на базата данни. При промяна на свързващия компонент се променя не само диалектът на заявката, но и целият програмен интерфейс за връзка, подготвяне, свързване (свързване) на параметри и четене на резултати.

Този подход е особено проблематичен, когато един и същ домейн модел трябва да бъде свързан последователно с повече от един тип база данни. В такъв случай приложението започва да натрупва условни пътища, специфични за дадения свързващ компонент SQL низове и дублирана логика за попълване на резултати. След определен размер на системата това намалява не само удобството, но и проверимостта на коректността.

Силната страна на този подход е неговата честност. Той не обещава повече абстракция, отколкото реално предоставя. Именно затова той често служи като базова линия, спрямо която могат да бъдат оценявани по-високите слоеве на абстракция.

\subsection{Асинхронност в съществуващите клиентски стекове}

Съществена част от съвременния контекст е, че платформено-специфичните клиентски библиотеки невинаги са само синхронни. \code{libpq} предоставя автомат с крайни състояния за асинхронно изпращане и получаване на заявки~\cite{libpq_async}, MySQL Конектор/C предлага двойки от тип \code{\_start}/\code{\_cont} за неблокиращи операции~\cite{mysql_capi_async}, \code{hiredis} има слой за функции за обратно извикване~\cite{hiredis}, а драйверът за Cassandra използва механизъм с бъдеща стойност и функция за обратно извикване~\cite{cassandra_futures}. Тези възможности обаче са силно специфични за съответния компонент както по интерфейс, така и по семантика на изпълнение.

Оттук следва важен извод: наличието на платформено-специфичен асинхронен програмен интерфейс само по себе си не решава проблема за обща C++ абстракция. Във всеки отделен случай приложението трябва да знае дали работи с цикъл за готовност на мрежово гнездо (готовност на мрежово гнездо), с регистрация на функция за обратно извикване, с периодично допитване до бъдеща стойност (периодично допитване до бъдеща стойност) или с вътрешно управлявани драйверни нишки. При MongoDB акцентът остава върху поддържането на пул и безопасна многонишкова употреба на \code{libmongoc}~\cite{mongoc_client_pool}, а при \code{libneo4j-client} преобладава синхронният модел заявка/отговор~\cite{libneo4j_client}. Именно тази хетерогенност прави необходим междинен асинхронен слой, който да бъде достатъчно честен към различията в свързващите компоненти и едновременно достатъчно единен за потребителския програмен интерфейс.

\subsection{Конструктори на заявки по време на изпълнение}

Библиотеки като \code{SOCI}~\cite{soci} и част от по-лекия инструментариум за изграждане на SQL търсят баланс между четимост и контрол. Те предлагат по-удобен C++ интерфейс за изграждане на заявки, но в много случаи продължават да разчитат на описване по време на изпълнение и последващо рендиране към SQL. В практиката това означава, че известна част от структурата на заявката е по-добре капсулирана, но смяната на модел за съхранение, независима от свързващия компонент, остава ограничена.

За настоящата разработка тези библиотеки са важни като отправна точка. Те показват, че изразителните програмни интерфейси (изразителен API) и типизирани елементи са полезни, но не са достатъчни сами по себе си. Същинският въпрос е дали компилаторът може да валидира цялата логическа структура на заявката и дали същата абстракция може да бъде пренесена извън SQL света.

\subsection{Типизирани SQL решения}

\code{sqlpp11}~\cite{sqlpp11} е особено интересен случай, защото стои между традиционния конструктор на заявки по време на изпълнение и по-строго статичния дизайн. Подобни библиотеки използват шаблони и типове, за да направят част от SQL структурата видима за компилатора. Това е съществена стъпка напред спрямо чистите конструктори по време на изпълнение, защото част от структурните грешки могат да бъдат хванати по-рано.

Същевременно обаче типизираното SQL решение обикновено остава дълбоко свързано с релационния модел. Дори когато е много силно като C++ програмен интерфейс, то най-често предполага свят, в който таблици, колони, съединявания (\code{JOIN}) и агрегати остават централната метафора. Това прави такива решения изключително ценни за SQL домейна, но по-трудно преносими към документни, графови или ключ-стойност системи.

Именно тук се очертава една от основните разлики с настоящата разработка. Тя не цели просто структуриране на SQL по време на компилация. Тя цели логически слой по време на компилация, който може да бъде материализиран различно според семейството на свързващия компонент.

\subsection{Класически рамки за обектно-релационно съпоставяне и кодогенерация}

\code{ODB}~\cite{odb} е характерен пример за рамка за обектно-релационно съпоставяне, която предлага богато картографиране между класове и релационни таблици, но използва генерация на код и външен инструментариум. Това улеснява потребителя на ниво употреба, но измества сложността към процеса на изграждане (процеса на изграждане) и към автоматично генерираните артефакти. Получава се силна зависимост между модела на данните, допълнителната стъпка в инструментариума и конкретните релационни компоненти.

В класическите ОРС рамки често се наблюдава и друго напрежение: колкото по-високо е удобството на ниво приложен код, толкова по-непрозрачен става пътят до реалната заявка и до експлоатационните характеристики на системата. Това не е задължително фатален проблем, но в академичния контекст на настоящата работа е важно ограничение, защото затруднява проследимостта между абстракцията и материализацията.

\subsection{Граници на релационно центрираните абстракции}

За приложения, които трябва да работят с документни, графови или ключ-стойност системи, класическият ОРС модел показва естествени ограничения. Част от концепциите --- като \code{JOIN}, външни ключове и строго таблична схема --- нямат директен еквивалент във всички модели за съхранение. Поради това обобщението на обектно-релационното съпоставяне отвъд релационния свят не може да бъде постигнато само чрез генериране на SQL.

Това наблюдение е важно и за специфичния програмен интерфейс, и за типизираните SQL решения. Дори когато те са много силни технически, обикновено предполагат, че SQL е централният език на света на данните. Настоящата разработка приема по-различна отправна точка: логическият модел трябва да бъде първичен, а релационната материализация --- само една от възможните интерпретации.

\begin{table}[H]
\centering
\caption{Ограничения на различните семейства решения спрямо целта за множество свързващи компоненти}
\label{tab:existing_solution_limitations}
\begin{tabularx}{\textwidth}{L{3.2cm}L{3.3cm}Y}
\toprule
\textbf{Семейство решения} & \textbf{Какво прави добре} & \textbf{Къде се появява ограничението} \\
\midrule
Специфичен интерфейс & максимален контрол и близост до драйвера & липса на общ логически модел и голямо дублиране при повече от един свързващ компонент \\
Конструктори по време на изпълнение & по-добра четимост и капсулиране на заявката & валидацията остава предимно по време на изпълнение и често е центрирана около SQL \\
Типизирани SQL решения & по-силна статична структура за SQL заявки & логиката обикновено остава ограничена в рамките на релационната метафора \\
Класическо ОРС & удобно съпоставяне на обекти за релационни системи & зависимост от кодогенерация и трудна преносимост към нерелационни модели \\
\bottomrule
\end{tabularx}
\end{table}

\subsection{Позициониране на настоящата разработка}

Предложената библиотека се позиционира като решение за обектно-релационно съпоставяне по време на компилация, което комбинира силна статична типизация, липса на външна кодогенерация, модел с конектори за пренасяне на логическата абстракция към различни целеви системи и базиран на корутини асинхронен слой за изпълнение. В сравнение с подхода със специфични програмни интерфейси тя намалява дублирането на логиката и риска от късно открити грешки. В сравнение с решенията с конструктори на заявки по време на изпълнение премества структурната валидация към компилатора. В сравнение с класическите ОРС рамки тя избягва външни генератори и поставя ясен акцент върху договора на конектора. В сравнение със специфичните за компонента асинхронни интерфейси тя предлага единен модел с корутини, без да твърди, че всички драйвери имат еднакви неблокиращи възможности.

Ключовата разлика обаче е по-дълбока. Настоящата разработка не разглежда техниките по време на компилация като средство за ``по-умен SQL конструктор'', а като начин за формализиране на логическия модел на достъпа до данни. Аналогично, тя не разглежда асинхронността като декоративна обвивка, а като отделен архитектурен проблем със собствени ограничения, правила за жизнения цикъл и специфични за компонента пътища на изпълнение. Точно това позволява следващите глави да преминат от теорията за корутините към архитектурата по време на компилация, асинхронния слой, договора на конектора и реалните сценарии с множество целеви системи.

\subsection{Изводи от сравнителния анализ}

Сравнението със съществуващите решения води до четири основни извода. Първо, подходът с платформено-специфични интерфейси остава незаменим като базова линия за контрол и производителност, но не предлага удовлетворително ниво на логическа абстракция. Второ, конструкторите на заявки по време на изпълнение и типизираните SQL решения подобряват работата със Structured Заявка Language (SQL), но не решават напълно проблема за логически слой, независим от компонента. Трето, класическите ОРС рамки предлагат висока абстракция, но често за сметка на сложност на инструментариума и ограничена пригодност извън релационния свят. Четвърто, платформено-специфичните асинхронни клиентски интерфейси са ценни, но показват толкова различни модели на изпълнение, че сами по себе си не осигуряват единна потребителска абстракция.

Точно това позициониране оправдава по-детайлното разглеждане в следващата глава на отделен, но ключов въпрос: кои езикови и теоретични механизми правят възможна базираната на корутини асинхронна архитектура, върху която по-късно стъпват ОРС моделът по време на компилация и изпълнителният слой за множество свързващи компоненти.
